% !TEX root = ../../main.tex
% File: part2/chapters4/chap4_4.tex

\section{Recipe 4: Fine-tuning cho Tác vụ Nhận dạng Thực thể (NER)}
\label{sec:recipe_ner_finetuning}

\textbf{Mục tiêu:} Fine-tune một mô hình Encoder-Only (ví dụ: RoBERTa) cho tác vụ NER. Recipe này nhấn mạnh vào việc tiền xử lý dữ liệu phức tạp hơn cho các tác vụ gán nhãn chuỗi.

\textbf{Thành phần chính:}
\begin{itemize}
    \item \textbf{\texttt{transformers}:} Sử dụng \texttt{AutoModelForTokenClassification} và \texttt{Trainer}.
    \item \textbf{\texttt{datasets}:} Tải bộ dữ liệu NER.
    \item \textbf{\texttt{evaluate} (với \texttt{seqeval}):} Để tính toán F1-score ở cấp độ thực thể.
\end{itemize}

\textbf{Các bước thực hiện:}
\begin{enumerate}
    \item \textbf{Tải dữ liệu và tokenizer:} Tải bộ dữ liệu CoNLL-2003.
    \item \textbf{Tiền xử lý và Căn chỉnh Nhãn:} Đây là bước phức tạp. Khi một từ bị tách thành nhiều subword bởi tokenizer, chúng ta cần đảm bảo các nhãn được căn chỉnh một cách chính xác.
    \item \textbf{Tải mô hình:} Tải một mô hình \texttt{AutoModelForTokenClassification}.
    \item \textbf{Thiết lập \texttt{Trainer}:} Sử dụng hàm \texttt{compute\_metrics} với \texttt{seqeval}.
    \item \textbf{Huấn luyện:} Bắt đầu quá trình fine-tuning.
\end{enumerate}


\textbf{Mã nguồn (tập trung vào bước tiền xử lý):}

\begin{example}{Fine-tuning cho Token Classification (NER)}{ex:ner_finetuning}
    \begin{minted}{python}
    # ... (cài đặt và import tương tự các recipe trước) ...
    from datasets import load_dataset
    from transformers import AutoTokenizer, AutoModelForTokenClassification, Trainer, TrainingArguments, DataCollatorForTokenClassification
    import evaluate
    import numpy as np
    
    # Bước 1: Tải dữ liệu
    raw_datasets = load_dataset("conll2003")
    label_names = raw_datasets["train"].features["ner_tags"].feature.names
    
    model_checkpoint = "roberta-base"
    tokenizer = AutoTokenizer.from_pretrained(model_checkpoint)
    
    # Bước 2: Hàm tiền xử lý và căn chỉnh nhãn
    def tokenize_and_align_labels(examples):
        tokenized_inputs = tokenizer(examples["tokens"], truncation=True, is_split_into_words=True)
        labels = []
        for i, label in enumerate(examples["ner_tags"]):
            word_ids = tokenized_inputs.word_ids(batch_index=i)
            previous_word_idx = None
            label_ids = []
            for word_idx in word_ids:
                if word_idx is None: # Special tokens like [CLS], [SEP]
                    label_ids.append(-100)
                elif word_idx != previous_word_idx: # First token of a new word
                    label_ids.append(label[word_idx])
                else: # Subsequent tokens of the same word
                    label_ids.append(-100) # Chỉ gán nhãn cho subword đầu tiên
                previous_word_idx = word_idx
            labels.append(label_ids)
        tokenized_inputs["labels"] = labels
        return tokenized_inputs
    
    tokenized_datasets = raw_datasets.map(tokenize_and_align_labels, batched=True)
    
    # Bước 3 & 4 & 5: Tải mô hình, thiết lập Trainer và Huấn luyện
    data_collator = DataCollatorForTokenClassification(tokenizer=tokenizer)
    seqeval = evaluate.load("seqeval")
    
    def compute_metrics(p):
        predictions, labels = p
        predictions = np.argmax(predictions, axis=2)
        # (Phần decode và tính toán tương tự ví dụ seqeval ở Chương 2.2)
        # ...
        return {"precision": ..., "recall": ..., "f1": ...} # Kết quả từ seqeval
    
    model = AutoModelForTokenClassification.from_pretrained(
        model_checkpoint, num_labels=len(label_names),
    )
    
    args = TrainingArguments("roberta-ner", eval_strategy="epoch", learning_rate=2e-5)
    trainer = Trainer(
        model, args,
        train_dataset=tokenized_datasets["train"],
        eval_dataset=tokenized_datasets["validation"],
        data_collator=data_collator,
        processing_class=tokenizer,
        compute_metrics=compute_metrics
    )
    trainer.train()
    \end{minted}
\end{example}
\textbf{Kết quả mong đợi:} Một mô hình được fine-tune chuyên biệt cho việc trích xuất các thực thể như Tên người, Tổ chức, Địa điểm từ văn bản.
